\section{AMC 10 10B 2006}
        \begin{enumerate}
        	\item (\textit{AMC 10 10B 2006}) What is $(-1)^{1} + (-1)^{2} + ... + (-1)^{2006}$ ?


 $\mathrm{(A) \ } -2006\qquad \mathrm{(B) \ } -1\qquad \mathrm{(C) \ } 0\qquad \mathrm{(D) \ } 1\qquad \mathrm{(E) \ } 2006$

 \textbf{Answer: }C



\textbf{Solution 1}

Since $-1$ raised to an odd integer is $-1$ and $-1$ raised to an even integer exponent is $1$: 

 $(-1)^{1} + (-1)^{2} + ... + (-1)^{2006} = (-1) + (1) + ... + (-1)+(1) = \boxed{\textbf{(C) }0}.$




\textbf{Solution 2}

The Using the formula for the first $n$ terms of a geometric series, with $n = 2006$, $a = (-1)^{1} = -1$, and $r = -1$, the sum can also be obtained:

 $\frac{a(1-r^n)}{1-r} = \frac{-1(1 - (-1)^{2006})}{1 + 1} = \frac{-1(1 - 1)}{2} = \frac{0}{2} = \boxed{\textbf{(C) }0}$



	\item (\textit{AMC 10 10B 2006}) For real numbers $x$ and $y$, define $x \mathop{\spadesuit} y = (x+y)(x-y)$. What is $3 \mathop{\spadesuit} (4 \mathop{\spadesuit} 5)$?


 $\mathrm{(A) \ } -72\qquad \mathrm{(B) \ } -27\qquad \mathrm{(C) \ } -24\qquad \mathrm{(D) \ } 24\qquad \mathrm{(E) \ } 72$

 \textbf{Answer: }A



\textbf{Solution 1}

Since $x \spadesuit y = (x+y)(x-y)$: 

 $3 \spadesuit (4 \spadesuit 5) = 3 \spadesuit((4+5)(4-5)) =  3 \spadesuit (-9) = (3+(-9))(3-(-9)) = \boxed{\textbf{(A)}-72}$




\textbf{Solution 2}

From \href{https://artofproblemsolving.com/wiki/index.php?title=Difference\_of\_squares}{difference of squares}, we have $x \spadesuit y = (x+y)(x-y) = x^2 - y^2$

 So:

 $3 \spadesuit (4 \spadesuit 5) = 3 \spadesuit(4^2 - 5^2) = 3 \spadesuit (16 - 25) = 3 \spadesuit (-9)= (3^2) - (-9)^2 =  9 - 81 = \boxed{\textbf{(A)}-72}$



	\item (\textit{AMC 10 10B 2006}) A football game was played between two teams, the Cougars and the Panthers. The two teams scored a total of 34 points, and the Cougars won by a margin of 14 points. How many points did the Panthers score? 


 $\mathrm{(A) \ } 10\qquad \mathrm{(B) \ } 14\qquad \mathrm{(C) \ } 17\qquad \mathrm{(D) \ } 20\qquad \mathrm{(E) \ } 24$

 \textbf{Answer: }A



\textbf{Solution}

Let $x$ be the number of points scored by the Cougars, and $y$ be the number of points scored by the Panthers. The problem is asking for the value of $y$. 
 \begin{align*} x+y &= 34 \\ x-y &= 14 \\ 2x &= 48 \\ x &= 24 \\ \end{align*}
The answer is $\boxed{\textbf{(A) } 10}$




\textbf{Solution 2}

$c$ is the amount the Cougars scored and $p$ is the score for Panthers. Since the Cougars won by 14 points, $c = p + 14$. Using substitution, $2p + 14 = 34$,$2p = 20$, and then$p = 10$.

 
 \begin{align*} p &= \boxed{\textbf{(A) }10} \\ \end{align*}


 -- leafy

 


	\item (\textit{AMC 10 10B 2006}) Circles of diameter 1 inch and 3 inches have the same center. The smaller circle is painted red, and the portion outside the smaller circle and inside the larger circle is painted blue. What is the ratio of the blue-painted area to the red-painted area? 


 
                    \begin{figure}[H]
                        \centering
                        \includegraphics[width=0.5\textwidth]{Problem/AMC_10_10B_2006_4_Problem_Figure_1.png}
                    \end{figure}
                    
$\mathrm{(A) \ } 2\qquad \mathrm{(B) \ } 3\qquad \mathrm{(C) \ } 6\qquad \mathrm{(D) \ } 8\qquad \mathrm{(E) \ } 9$

 \textbf{Answer: }D



\textbf{Solution}

The area painted red is equal to the area of the smaller circle and the area painted blue is equal to the area of the larger circle minus the area of the smaller circle. 

 So we have:

 
 \begin{align*} A_{red}&=\pi\left(\frac{1}{2}\right)^2=\frac{\pi}{4}\\ A_{blue}&=\pi\left(\frac{3}{2}\right)^2-\pi\left(\frac{1}{2}\right)^2=2\pi\\ \end{align*}


 So the desired ratio is: 

 $\frac{A_{blue}}{A_{red}}=\frac{2\pi}{\frac{\pi}{4}}=2\cancel{\pi}\cdot \frac{4}{\cancel{\pi}}=\boxed{\textbf{(D) }8}.$



	\item (\textit{AMC 10 10B 2006}) A $2 \times 3$ rectangle and a $3 \times 4$ rectangle are contained within a square without overlapping at any point, and the sides of the square are parallel to the sides of the two given rectangles. What is the smallest possible area of the square? 


 $\mathrm{(A) \ } 16\qquad \mathrm{(B) \ } 25\qquad \mathrm{(C) \ } 36\qquad \mathrm{(D) \ } 49\qquad \mathrm{(E) \ } 64$

 \textbf{Answer: }B



\textbf{Solution 1}

By placing the $2 \times 3$ rectangle adjacent to the $3 \times 4$ rectangle with the 3 side of the $2 \times 3$ rectangle next to the 4 side of the $3 \times 4$ rectangle, we get a figure that can be completely enclosed in a square with a side length of 5. The area of this square is $5^2 = 25$.

 Since placing the two rectangles inside a $4 \times 4$ square must result in overlap, the smallest possible area of the square is $0$.

 So the answer is $\boxed{\textbf{(B) }25}$.

 



\textbf{Solution 2}

The area of a $2\times3$ rectangle and a $3\times4$ rectangle combined is $18$, so a $4\times4$ square is impossible without overlapping. Thus, the next smallest square is a $5\times5$, which works, so the answer is $\boxed{B)25}$.

 Note: If you do this, always check to see if it fits, because this doesn't always work. For example, a $3\times3$ and a $3\times4$ doesn't fit into a $5\times5$, even though their combined area is $21$.

 


	\item (\textit{AMC 10 10B 2006}) A region is bounded by semicircular arcs constructed on the side of a square whose sides measure $\frac{2}{\pi}$, as shown. What is the perimeter of this region? 


 
                    \begin{figure}[H]
                        \centering
                        \includegraphics[width=0.5\textwidth]{Problem/AMC_10_10B_2006_6_Problem_Figure_1.png}
                    \end{figure}
                    
$\mathrm{(A) \ } \frac{4}{\pi}\qquad \mathrm{(B) \ } 2\qquad \mathrm{(C) \ } \frac{8}{\pi}\qquad \mathrm{(D) \ } 4\qquad \mathrm{(E) \ } \frac{16}{\pi}$

 \textbf{Answer: }D



\textbf{Solution}

Since the side of the square is the diameter of the semicircle, the radius of the semicircle is $\frac{1}{2}\cdot\frac{2}{\pi}=\frac{1}{\pi}$.

 Since the length of one of the semicircular \href{https://artofproblemsolving.com/wiki/index.php?title=Arc}{arcs} is half the circumference of the corresponding circle, the length of one arc is $\frac{1}{2}\cdot2\cdot\pi\cdot\frac{1}{\pi}=1$. 

 Since the desired perimeter is made up of four of these arcs, the perimeter is $4\cdot1=\boxed{\textbf{(D) }4}$.

 


	\item (\textit{AMC 10 10B 2006}) Which of the following is equivalent to $\sqrt{\frac{x}{1-\frac{x-1}{x}}}$ when $x < 0$?


 $\mathrm{(A) \ } -x\qquad \mathrm{(B) \ } x\qquad \mathrm{(C) \ } 1\qquad \mathrm{(D) \ } \sqrt{\frac{x}{2}}\qquad \mathrm{(E) \ } x\sqrt{-1}$

 \textbf{Answer: }A



\textbf{Solution 1}

$\sqrt{\frac{x}{1-\frac{x-1}{x}}} =  \sqrt{\frac{x}{\frac{x}{x}-\frac{x-1}{x}}} = \sqrt{\frac{x}{\frac{x-(x-1)}{x}}} = \sqrt{\frac{x}{\frac{1}{x}}} = \sqrt{x^2} = |x|$

 Since $x<0,|x|= \boxed{\textbf{(A)}-x}$




\textbf{Solution 2}

To confirm the answer, inputting a negative value into $x$ may help. For the ease of computation, if $x=-3$, $\sqrt{\frac{-3}{1-\frac{4}{3}}}=\sqrt{\frac{-3}{\frac{-1}{3}}}=\sqrt{9}=3$. As no other option choice fits, $\boxed{\textbf{(A)}-x}$ is the correct solution. Note: If $x=-1$ was chosen, C would have fit as well. Make sure to avoid making this mistake.

 


	\item (\textit{AMC 10 10B 2006}) A square of area 40 is inscribed in a semicircle as shown. What is the area of the semicircle? 


 
                    \begin{figure}[H]
                        \centering
                        \includegraphics[width=0.5\textwidth]{Problem/AMC_10_10B_2006_8_Problem_Figure_1.png}
                    \end{figure}
                    
$\mathrm{(A) \ } 20\pi\qquad \mathrm{(B) \ } 25\pi\qquad \mathrm{(C) \ } 30\pi\qquad \mathrm{(D) \ } 40\pi\qquad \mathrm{(E) \ } 50\pi$

 \textbf{Answer: }B



\textbf{Solution}

Since the area of the square is $40$, the length of a side is $\sqrt{40}=2\sqrt{10}$. The distance between the center of the semicircle and one of the bottom vertices of the square is half the length of the side, which is $\sqrt{10}$.

 Using the \href{https://artofproblemsolving.com/wiki/index.php?title=Pythagorean\_Theorem}{Pythagorean Theorem} to find the radius $r$ of the semicircle, $r^2 = (2\sqrt{10})^2 + (\sqrt{10})^2 = 50$. So, the area of the semicircle is $\frac{1}{2}\cdot \pi \cdot 50 = \boxed{\textbf{(B) }25\pi}$.

 


	\item (\textit{AMC 10 10B 2006}) Francesca uses 100 grams of lemon juice, 100 grams of sugar, and 400 grams of water to make lemonade. There are 25 calories in 100 grams of lemon juice and 386 calories in 100 grams of sugar. Water contains no calories. How many calories are in 200 grams of her lemonade? 


 $\mathrm{(A) \ } 129\qquad \mathrm{(B) \ } 137\qquad \mathrm{(C) \ } 174\qquad \mathrm{(D) \ } 233\qquad \mathrm{(E) \ } 411$

 \textbf{Answer: }B



\textbf{Solution}

The calorie to gram ratio of Francesca's lemonade is $\frac{25+386+0}{100+100+400}=\frac{411\textrm{ calories}}{600\textrm{ grams}}=\frac{137\textrm{ calories}}{200\textrm{ grams}}$ 

 So in $200\textrm{ grams}$ of Francesca's lemonade there are $200\textrm{ grams}\cdot\frac{137\textrm{ calories}}{200\textrm{ grams}}=\boxed{\textbf{(B) }137\textrm{ calories}}.$



	\item (\textit{AMC 10 10B 2006}) In a triangle with integer side lengths, one side is three times as long as a second side, and the length of the third side is 15. What is the greatest possible perimeter of the triangle? 


 $\mathrm{(A) \ } 43\qquad \mathrm{(B) \ } 44\qquad \mathrm{(C) \ } 45\qquad \mathrm{(D) \ } 46\qquad \mathrm{(E) \ } 47$

 \textbf{Answer: }A



\textbf{Solution}

Let $x$ be the length of the first side. 

 The lengths of the sides are: $x$, $3x$, and $15$. 

 By the \href{https://artofproblemsolving.com/wiki/index.php?title=Triangle\_Inequality}{Triangle Inequality},

 $3x < x + 15$
$2x < 15$
$x < \frac{15}{2}$

 The greatest integer satisfying this inequality is $7$.

 So the greatest possible perimeter is $7 + 3\cdot7 + 15 =\boxed{\textbf{(A) } 43}$



	\item (\textit{AMC 10 10B 2006}) What is the tens digit in the sum $7!+8!+9!+...+2006!$
$\mathrm{(A) \ } 1\qquad \mathrm{(B) \ } 3\qquad \mathrm{(C) \ } 4\qquad \mathrm{(D) \ } 6\qquad \mathrm{(E) \ } 9$

 \textbf{Answer: }C



\textbf{Solution}

Since $10!$ is divisible by $100$, any \href{https://artofproblemsolving.com/wiki/index.php?title=Factorial}{factorial} greater than $10!$ is also divisible by $100$. The last two \href{https://artofproblemsolving.com/wiki/index.php?title=Digit}{digits} of all factorials greater than $10!$ are $00$, so the last two digits of $10!+11!+...+2006!$ are $00$. (*)

 So all that is needed is the tens digit of the sum $7!+8!+9!$
$7!+8!+9!=5040+40320+362880=408240$

 So the tens digit is $\boxed{\textbf{(C) }4}$.

 (*) A slightly faster method would be to take the $\pmod {100}$ residue of $7! + 8! + 9!.$ Since $7! = 5040,$ we can rewrite the sum as 
 \[5040 + 8\cdot 5040 + 72\cdot 5040 \equiv 40 + 8\cdot 40 + 72\cdot 40 = 40 + 320 + 2880 \equiv 40 \pmod{100}.\] Since the last two digits of the sum is $40$, the tens digit is $\boxed{\textbf{(C) }4}$.

 


	\item (\textit{AMC 10 10B 2006}) The lines $x=\frac{1}{4}y+a$ and $y=\frac{1}{4}x+b$ intersect at the point $(1,2)$. What is $a+b$?


 $\mathrm{(A) \ } 0\qquad \mathrm{(B) \ } \frac{3}{4}\qquad \mathrm{(C) \ } 1\qquad \mathrm{(D) \ } 2\qquad \mathrm{(E) \ } \frac{9}{4}$

 \textbf{Answer: }E



\textbf{Solution 1}

Since $(1,2)$ is a solution to both equations, plugging in $x=1$ and $y=2$ will give the values of $a$ and $b$. 

 $1 = \frac{1}{4} \cdot 2 + a$
$a = \frac{1}{2}$
$2 = \frac{1}{4} \cdot 1 + b$
$b = \frac{7}{4}$

 So $a+b = \frac{1}{2} + \frac{7}{4} = \boxed{\textbf{(E) }\frac{9}{4}}$.

 



\textbf{Solution 2}

Substituting $x=1$ and $y=2$ into the equations gives $1=\frac{2}{4}+a\quad\text{and}\quad 2=\frac{1}{4}+b.$ It follows that $a+b=\left(1-\frac{2}{4}\right)+\left(2-\frac{1}{4}\right)=3 - \frac{3}{4}=\boxed{\frac{9}{4}}.$




\textbf{Solution 3}

Because $a=x- \frac{y}{4}\quad\text{and}\quad b=y-\frac{x}{4} \quad\text{we have}\quad a +b = \frac{3}{4}(x+y).$ Since $x=1$ when $y=2$, this implies that $a+b = \frac{3}{4}(1 + 2) = \boxed{\frac{9}{4}}$.

 


	\item (\textit{AMC 10 10B 2006}) Joe and JoAnn each bought 12 ounces of coffee in a 16 ounce cup. Joe drank 2 ounces of his coffee and then added 2 ounces of cream. JoAnn added 2 ounces of cream, stirred the coffee well, and then drank 2 ounces. What is the resulting ratio of the amount of cream in Joe's coffee to that in JoAnn's coffee? 


 $\mathrm{(A) \ } \frac{6}{7}\qquad \mathrm{(B) \ } \frac{13}{14}\qquad \mathrm{(C) \ }1 \qquad \mathrm{(D) \ } \frac{14}{13}\qquad \mathrm{(E) \ } \frac{7}{6}$

 \textbf{Answer: }E



\textbf{Solution}

After drinking and adding cream, Joe's cup has $2$ ounces of cream. 

 After adding cream to her cup, JoAnn's cup had $14$ ounces of liquid. By stirring and then drinking $2$ ounces out of the $14$ ounces of liquid, she drank $\frac{2}{14}=\frac{1}{7}$th of the cream. So there are $2\cdot\frac{6}{7}=\frac{12}{7}$ ounces of cream left. 

 So the desired ratio is: $2 \div \frac{12}{7}= \boxed{\textbf{(E) }\frac{7}{6}}$.

 


	\item (\textit{AMC 10 10B 2006}) Let $a$ and $b$ be the roots of the equation $x^2-mx+2=0$. Suppose that $a+\frac{1}{b}$ and $b+\frac{1}{a}$ are the roots of the equation $x^2-px+q=0$. What is $q$?


 $\mathrm{(A) \ } \frac{5}{2}\qquad \mathrm{(B) \ } \frac{7}{2}\qquad \mathrm{(C) \ } 4\qquad \mathrm{(D) \ } \frac{9}{2}\qquad \mathrm{(E) \ } 8$

 \textbf{Answer: }D



\textbf{Solution}

In a \href{https://artofproblemsolving.com/wiki/index.php?title=Quadratic\_equation}{quadratic equation} of the form $x^2 + bx + c = 0$, the product of the \href{https://artofproblemsolving.com/wiki/index.php?title=Root}{roots} is $c$ (\href{https://artofproblemsolving.com/wiki/index.php?title=Vieta\%27s\_Formulas}{Vieta's Formulas}).

 Using this property, we have that $ab=2$ and

 $q = (a+\frac{1}{b})\cdot(b+\frac{1}{a}) = \frac{ab+1}{b} \cdot \frac{ab+1}{a}  = \frac{(ab+1)^2}{ab} = \frac{(2+1)^2}{2} = \boxed{\textbf{(D) }\frac{9}{2}}$.

 
            \begin{itemize}
        \item Notice the fact that we never actually found the roots.
            \end{itemize}
            



\textbf{Solution 2}

Assume \href{https://artofproblemsolving.com/wiki/index.php?title=Without\_loss\_of\_generality}{without loss of generality} that $m=3$. We can factor the equation $x^2-3x+2=0$ into $(x-1)(x-2)=0$. Therefore, $a=1$ and $b=2$. Using these values, we find $a+\frac1b=\frac32$ and $b+\frac1a=3$. By \href{https://artofproblemsolving.com/wiki/index.php?title=Vieta\%27s\_formulas}{Vieta's formulas}, $q$ is the product of the roots of $x^2-px+q=0$, which are $a+\frac1b$ and $b+\frac1a$. Therefore, $q=\left(a+\frac1b\right)\left(b+\frac1a\right)=\frac32\cdot3=\boxed{\textbf{(D) }\frac92}$.

 


	\item (\textit{AMC 10 10B 2006}) Rhombus $ABCD$ is similar to rhombus $BFDE$. The area of rhombus $ABCD$ is $24$ and $\angle BAD = 60^\circ$. What is the area of rhombus $BFDE$? 


 
                    \begin{figure}[H]
                        \centering
                        \includegraphics[width=0.5\textwidth]{Problem/AMC_10_10B_2006_15_Problem_Figure_1.png}
                    \end{figure}
                    
$\mathrm{(A) \ } 6\qquad \mathrm{(B) \ } 4\sqrt{3}\qquad \mathrm{(C) \ } 8\qquad \mathrm{(D) \ } 9\qquad \mathrm{(E) \ } 6\sqrt{3}$

 \textbf{Answer: }C



\textbf{Solution 1}

Using the property that opposite angles are equal in a \href{https://artofproblemsolving.com/wiki/index.php?title=Rhombus}{rhombus}, $\angle DAB = \angle DCB = 60 ^\circ$ and $\angle ADC = \angle ABC = 120 ^\circ$. It is easy to see that rhombus $ABCD$ is made up of \href{https://artofproblemsolving.com/wiki/index.php?title=Equilateral\_triangle}{equilateral triangles} $DAB$ and $DCB$. Let the lengths of the sides of rhombus $ABCD$ be $s$. 

 The longer \href{https://artofproblemsolving.com/wiki/index.php?title=Diagonal}{diagonal} of rhombus $BFDE$ is $BD$. Since $BD$ is a side of an equilateral triangle with a side length of $s$, $BD = s$. The longer diagonal of rhombus $ABCD$ is $AC$. Since $AC$ is twice the length of an altitude of of an equilateral triangle with a side length of $s$, $AC = 2 \cdot \frac{s\sqrt{3}}{2} = s\sqrt{3}$.

 The ratio of the longer diagonal of rhombus $BFDE$ to rhombus $ABCD$ is $\frac{s}{s\sqrt{3}} = \frac{\sqrt{3}}{3}$. Therefore, the ratio of the \href{https://artofproblemsolving.com/wiki/index.php?title=Area}{area} of rhombus $BFDE$ to rhombus $ABCD$ is $\left( \frac{\sqrt{3}}{3} \right) ^2 = \frac{1}{3}$.

 Let $x$ be the area of rhombus $BFDE$. Then $\frac{x}{24} = \frac{1}{3}$, so $x = \boxed{\textbf{(C) }8}$.

 



\textbf{Solution 2}

Triangle DAB is equilateral so triangles $DEA$, $AEB$, $BED$, $BFD$, $BFC$ and $CFD$ are all congruent with angles $30^\circ$, $30^\circ$ and $120^\circ$ from which it follows that rhombus $BFDE$ has one third the area of rhombus $ABCD$ i.e. $8 \Longrightarrow \boxed{\textbf{(C) }8}$.Note: A quick way to visualize this method is to draw the line $DB$ and notice the two equilateral triangles $\triangle ADB$ and $\triangle DBC$.

 



\textbf{Solution 3 (A bit more math involved)}

We can extend line $DE$, meeting line $AB$ at $G$. Similarly, we can extend line $DF$ to meet line $BC$ at $H$. We can see with some simple math that triangle $ADG$ is a $30$-$60$-$90$ triangle, so we can call line $AG$ as $x$, line $DG$ as $x\sqrt{3}$, and line $AD$ as $2x$ (because of the $30$-$60$-$90$ triangle side proportions). 

 We can also see that line $AD$ is a base of rhombus $ABCD$, and line $DH$ is a height. Since triangle $DHC$ is also a $30$-$60$-$90$ triangle, line $DH$ is also $x\sqrt{3}$. Since the question told us that the area of rhombus $ABCD$ is $24$, we can make the following equation:

 $2x \cdot x\sqrt{3} = 24$

 Solving for x:

 $2x^2\sqrt{3} = 24$
$x^2\sqrt{3} = 12$
$x^2 = \frac{12}{\sqrt{3}}$
$x^2 = 4\sqrt{3}$
$x = 2\sqrt{\sqrt{3}}$

 
Since the question is to find the area of rhombus $BFDE$, to find the answer, we can just multiply base $DE$ with the rhombus's height. We'll start by finding the height: instantly we can see that $GB$ is the height. Since all the sides of a rhombus are equal, and we found earlier that the side length is $2x$, if $AG$ is $x$, that means $GB$ is the same length as $AG$ - that is to say, 

 $GB = AG = 2\sqrt{\sqrt{3}}$

 
Now to find the base. We can see that to find the base, we can simply just subtract the length of line $EG$ from the length of line $DG$. Since $DG$ is $x\sqrt{3}$, and $x$ is $2\sqrt{\sqrt{3}}$, that makes

 $DG = 2\sqrt{\sqrt{3}} \cdot \sqrt{3} = 2\sqrt{3\sqrt{3}}$ 

 
Now to find $EG$: We can see with simple math that triangle $EGB$ is also a $30$-$60$-$90$ triangle, which means that $EG = \frac{GB}{\sqrt{3}}$. Previously, we found out that $GB$ is $2\sqrt{\sqrt{3}}$, so:

 $EG = \frac{2\sqrt{\sqrt{3}}}{\sqrt{3}} = \frac{2\sqrt{3\sqrt{3}}}{3}$

 
Now we can find the base:

 $DG - EG = 2\sqrt{3\sqrt{3}} - \frac{2\sqrt{3\sqrt{3}}}{3} = \frac{4\sqrt{3\sqrt{3}}}{3}$

 
Multiplying the newly found base by the height we found earlier:

 $\frac{4\sqrt{3\sqrt{3}}}{3} \cdot 2\sqrt{\sqrt{3}} = \frac{8\sqrt{9}}{3} = \frac{24}{3} = \boxed{\textbf{(C) }8}$



	\item (\textit{AMC 10 10B 2006}) Leap Day, February 29, 2004, occurred on a Sunday. On what day of the week will Leap Day, February 29, 2020, occur? 


 $\mathrm{(A) \ } \textrm{Tuesday} \qquad \mathrm{(B) \ } \textrm{Wednesday} \qquad \mathrm{(C) \ } \textrm{Thursday} \qquad \mathrm{(D) \ } \textrm{Friday} \qquad \mathrm{(E) \ } \textrm{Saturday}$

 \textbf{Answer: }E



\textbf{Solution}

There are $365$ days in a year, plus $1$ extra day if there is a Leap Day, which occurs on years that are multiples of $4$ (with a few exceptions that don't affect this problem).  

 Therefore, the number of days between Leap Day $2004$ and Leap Day $2020$ is:  

 $16 \cdot 365 + 4 \cdot 1 = 5844$

 Since the days of the week repeat every $7$ days and $5844 \equiv -1 \bmod{7}$, the day of the week Leap Day $2020$ occurs is the day of the week the day before Leap Day $2004$ occurs, which is $\boxed{\textbf{(E) }\textrm{Saturday}}$.

 



\textbf{Solution 2 (Feasible Shortcuts)}

Since every non-leap year there is one day extra after the $52$ weeks, we can deduce that if we were to travel forward $x$ amount of non-leap years, then the answer would be "Sunday + $x$" days.

 However, we also have leap-years in our pool of years traveled forward, and for every leap-year that passes, we have two extra days after the 52 weeks. So we would travel "Sunday + $2x$" days. 

 Mixing these two together, we get $4$ leap years $(2008, 2012, 2016, 2020)$, and $12$ non-leap years. We thus get $12 + 2\cdot 4 = 20$ "extra days" after Sunday. Since seven days are in a week, we can get $20 \bmod{7} \equiv 6$, and six days after Sunday is $\boxed{\textbf{(E) }\textrm{Saturday}}$.

 
            \begin{itemize}
        \item Note how we did not require any complex or time-consuming computation, and thus it will save crucial time, especially during a test like the AMC 10.
            \end{itemize}
            


	\item (\textit{AMC 10 10B 2006}) Bob and Alice each have a bag that contains one ball of each of the colors blue, green, orange, red, and violet. Alice randomly selects one ball from her bag and puts it into Bob's bag. Bob then randomly selects one ball from his bag and puts it into Alice's bag. What is the probability that after this process the contents of the two bags are the same? 


 $\mathrm{(A) \ } \frac{1}{10}\qquad \mathrm{(B) \ } \frac{1}{6}\qquad \mathrm{(C) \ } \frac{1}{5}\qquad \mathrm{(D) \ } \frac{1}{3}\qquad \mathrm{(E) \ } \frac{1}{2}$

 \textbf{Answer: }D



\textbf{Solution}

Since there are the same amount of total balls in Alice's bag as in Bob's bag, and there is an equal chance of each ball being selected, the color of the ball that Alice puts in Bob's bag doesn't matter. \href{https://artofproblemsolving.com/wiki/index.php?title=Without\_loss\_of\_generality}{Without loss of generality}, let the ball Alice puts in Bob's bag be red. 

 For both bags to have the same contents, Bob must select one of the $2$ red balls out of the $6$ balls in his bag. 

 So the desired probability is $\frac{2}{6} = \boxed{\textbf{(D) }\frac{1}{3}}$.

 


	\item (\textit{AMC 10 10B 2006}) Let $a_1 , a_2 , ...$ be a sequence for which


 $a_1=2$ , $a_2=3$, and $a_n=\frac{a_{n-1}}{a_{n-2}}$ for each positive integer $n \ge 3$. 


 What is $a_{2006}$?


 $\mathrm{(A) \ } \frac{1}{2}\qquad \mathrm{(B) \ } \frac{2}{3}\qquad \mathrm{(C) \ } \frac{3}{2}\qquad \mathrm{(D) \ } 2\qquad \mathrm{(E) \ } 3$

 \textbf{Answer: }E



\textbf{Solution 1 (Kinda Cheese)}

Looking at the first few terms of the sequence:

 $a_1=2 , a_2=3 , a_3=\frac{3}{2}, a_4=\frac{1}{2} , a_5=\frac{1}{3} , a_6=\frac{2}{3} , a_7=2 , a_8=3 , ....$

 Clearly, the sequence repeats every 6 terms. 

 Since $2006 \equiv 2\bmod{6}$,

 $a_{2006} = a_2 = \boxed{\textbf{(E) }3}$




\textbf{Solution 2}

$a_n = \frac{a_{n-1}}{a_{n-2}} = \frac{\frac{a_{n-2}}{a_{n-3}}}{a_{n-2}} = \frac{1}{a_{n-3}}$ , so $a_n = a_{n-6}$ and because $2006 = 2 + 334 \times 6$ , so $a_{2006} = a_2 = \boxed{\textbf{(E) }3}$



	\item (\textit{AMC 10 10B 2006}) A circle of radius $2$ is centered at $O$. Square $OABC$ has side length $1$. Sides $AB$ and $CB$ are extended past $B$ to meet the circle at $D$ and $E$, respectively. What is the area of the shaded region in the figure, which is bounded by $BD$, $BE$, and the minor arc connecting $D$ and $E$?


 
                    \begin{figure}[H]
                        \centering
                        \includegraphics[width=0.5\textwidth]{Problem/AMC_10_10B_2006_19_Problem_Figure_1.png}
                    \end{figure}
                    
$\mathrm{(A) \ } \frac{\pi}{3}+1-\sqrt{3}\qquad \mathrm{(B) \ } \frac{\pi}{2}(2-\sqrt{3})\qquad \mathrm{(C) \ } \pi(2-\sqrt{3})\qquad \mathrm{(D) \ } \frac{\pi}{6}+\frac{\sqrt{3}+1}{2}\qquad \mathrm{(E) \ } \frac{\pi}{3}-1+\sqrt{3}$

 \textbf{Answer: }A



\textbf{Solution 1}

The shaded area is equivalent to the area of sector $DOE$ minus the area of triangle $DOE$ plus the area of triangle $DBE$. 

 Using the \href{https://artofproblemsolving.com/wiki/index.php?title=Pythagorean\_Theorem}{Pythagorean Theorem}, $(DA)^2=(CE)^2=2^2-1^2=3$ so $DA=CE=\sqrt{3}$.

 Clearly, $DOA$ and $EOC$ are $30-60-90$ triangles with $\angle EOC = \angle DOA = 60^\circ$. Since $OABC$ is a square, $\angle COA = 90^\circ$.

 $\angle DOE$ can be found by doing some subtraction of angles. 

 $\angle COA - \angle DOA = \angle DOC = \angle EOA$
$90^\circ - 60^\circ = \angle EOA = 30^\circ$
$\angle DOA - \angle EOA = \angle DOE$
$60^\circ - 30^\circ = \angle DOE = 30^\circ$

 So, the area of sector $DOE$ is $\frac{30}{360} \cdot \pi \cdot 2^2 = \frac{\pi}{3}$.

 The area of triangle $DOE$ is $\frac{1}{2}\cdot 2 \cdot 2 \cdot \sin 30^\circ = 1$.

 Since $AB=CB=1$ , $DB=EB=(\sqrt{3}-1)$. So, the area of triangle $DBE$ is $\frac{1}{2} \cdot (\sqrt{3}-1)^2 = 2-\sqrt{3}$. Therefore, the shaded area is $(\frac{\pi}{3}) - (1) + (2-\sqrt{3}) = \boxed{\textbf{(A) }\frac{\pi}{3} + 1 - \sqrt{3}}$

 OR

 $\triangle{ODA}$ has the same height as $\triangle{OBD}$ which is $1.$

 We already know that $BD = \sqrt{3} - 1.$

 Therefore the area of $\triangle{OBD}$ is $(\sqrt{3}-1) \cdot 1 \cdot \frac{1}{2} = \frac{\sqrt{3}-1}{2}.$

 Since $\triangle{OBD} = \triangle{OBE} = \frac{\sqrt{3}-1}{2}.$

 Therefore the sum of the areas is $2 \cdot \frac{\sqrt{3}-1}{2} = \sqrt{3}-1.$

 Then the area of the shaded area becomes $\frac{\pi}{3} - (\sqrt{3} - 1) = \boxed{\textbf{(A) }\frac{\pi}{3} +1 - \sqrt{3}}.$




\textbf{Solution 2}

From the Pythagorean Theorem, we can see that $DA$ is $\sqrt{3}$. Then, $DB = DA - BA = \sqrt{3} - 1$. The area of the shaded element is the area of sector $DOE$ minus the areas of triangle $DBO$ and triangle $EBO$ combined. Below is an image to help.

 
                    \begin{figure}[H]
                        \centering
                        \includegraphics[width=0.5\textwidth]{Solution/AMC_10_10B_2006_19_Solution_Figure_1.png}
                    \end{figure}
                    

 Using the Base Altitude formula, where $DB$ and $BE$ are the bases and $OA$ and $CO$ are the altitudes, respectively, $[DBO] = [EBO] = \frac{\sqrt{3}-1}{2}$. The area of sector $DOE$ is $\frac{1}{12}$ of circle $O$. The area of circle $O$ is $4\pi$, and therefore we have the area of sector $DBE$ to be $\boxed{\textbf{(A) }\frac{\pi}{3} + 1 - \sqrt{3}}$.

 



\textbf{Solution 3 (Using Answer Choices)}

Like the first solutions, you find that the area of sector $DOE$ is $\frac{\pi}{3}$. We also know that the triangles will not be in terms of ${\pi}$. Looking at the answers, choices $\text{(A)}$ and $\text{(E)}$ both contain $\frac{\pi}{3}$. However, based on the diagram, we observe that the answer must be less than $\frac {\pi}{3}$. Only $\boxed{\textbf{(A) }\frac{\pi}{3} + 1 - \sqrt{3}}$ consists of a value less than $\frac{\pi}{3}$.

 



\textbf{Solution 4}

We calculate the shaded area by subtracting the unshaded area in the quarter circle from a quarter circle, which consists of two $30 ^{\circ}$ sectors and two 30-60-90 triangles minus a square. The area of the shaded region is  $\frac {\pi \times 2^2}{4} - 2 ( \frac {\pi \times 2^2}{12} + \frac {\sqrt 3}{2} - 1 ) = \boxed{\textbf{(A) }\frac{\pi}{3} + 1 - \sqrt{3}}$



	\item (\textit{AMC 10 10B 2006}) In rectangle $ABCD$, we have $A=(6,-22)$, $B=(2006,178)$, $D=(8,y)$, for some integer $y$. What is the area of rectangle $ABCD$?


 $\mathrm{(A) \ } 4000\qquad \mathrm{(B) \ } 4040\qquad \mathrm{(C) \ } 4400\qquad \mathrm{(D) \ } 40,000\qquad \mathrm{(E) \ } 40,400$

 \textbf{Answer: }E



\textbf{Solution 1}

Let the slope of $AB$ be $m_1$ and the slope of $AD$ be $m_2$.

 $m_1 = \frac{178-(-22)}{2006-6} = \frac{1}{10}$
$m_2 = \frac{y-(-22)}{8-6} = \frac{y+22}{2}$

 Since $AB$ and $AD$ form a right angle: 

 $m_2 = -\frac{1}{m_1}$
$m_2 = -10$
$\frac{y+22}{2} = -10$
$y = -42$

 Using the \href{https://artofproblemsolving.com/wiki/index.php?title=Distance\_formula}{distance formula}:

 $AB = \sqrt{ (2006-6)^2 + (178-(-22))^2 } = \sqrt{ (2000)^2 + (200)^2 } = 200\sqrt{101}$
$AD = \sqrt{ (8-6)^2 + (-42-(-22))^2 } = \sqrt{ (2)^2 + (-20)^2 } = 2\sqrt{101}$

 Therefore the area of rectangle $ABCD$ is $200\sqrt{101}\cdot2\sqrt{101} = 40,400 \Rightarrow E$




\textbf{Solution 2}

This solution is the same as Solution 1 up to the point where we find that $y=-42$.

 We build right triangles so we can use the Pythagorean Theorem.  The triangle with hypotenuse $AB$ has legs $200$ and $2000$, while the triangle with hypotenuse $AD$ has legs $2$ and $20$.  Aha!  The two triangles are similar by SAS, with one triangle having side lengths $100$ times the other!

 Let $AD=x$.  Then from our reasoning above, we have $AB=100x$.  Finally, the area of the rectangle is $100x(x)=100x^2=100(20^2+2^2)=100(400+4)=100(404)=\boxed{40400 \text{  (E)}}$.

 



\textbf{Solution 3}

We do not need to solve for y. We form a right triangle with $AB$ as the hypotenuse and two adjacent sides lengths 200 and 2000, respectively. We form another right triangle with $AD$ as the hypotenuse and 2 is one of the lengths of the adjacent sides. Those two triangles are similar because $AD$ and $AB$ are perpendicular. $\frac{AB}{AD} = \frac{200}{2}$, so the area $AB \cdot AD = \frac {AB^2}{100} = \frac {2000^2 + 200^2}{100} = \boxed{40400 \text{  (E)}}$




\textbf{Solution 4 (answer choices)}

In order to find the area of rectangle $ABCD$, we need to find $AB$ first. Using the \href{https://artofproblemsolving.com/wiki/index.php?title=Distance\_formula}{distance formula}, we can derive:

 $AB = \sqrt{(2006-6)^2 + (178-(-22))^2} = \sqrt{(2000)^2 + (200)^2 } = 200\sqrt{101}$ 

 Now we can look at the answer choices. Because all of them are integers, then we know that $AD$ has to also contain $\sqrt{101}$ to ensure that $AB*AD$ is an integer. Once you know that, you can guarantee that the answer must be a multiple of 101. (If you're wondering why there can't be a fraction like $1/101$, there can't because $y$ is an integer the number under the square root in $AD$, is an integer.) Having that information narrows down the answer choices to just $(B)$ $4040$, and $(E)$ $40400$.

 Looking back at what we found for $AB$, which is $200\sqrt{101}$, and what we know about $AD$, that it has to contain $\sqrt{101}$. we know that the answer is at least $200\sqrt{101}*\sqrt{101}=20200$, which is already greater than answer choice $(B)$, which is $4040$. From this, we can conclude that the answer is $\boxed{\textbf{(E) }  40400}$



	\item (\textit{AMC 10 10B 2006}) For a particular peculiar pair of dice, the probabilities of rolling $1$, $2$, $3$, $4$, $5$, and $6$, on each die are in the ratio $1:2:3:4:5:6$. What is the probability of rolling a total of $7$ on the two dice? 


 $\mathrm{(A) \ } \frac{4}{63}\qquad \mathrm{(B) \ } \frac{1}{8}\qquad \mathrm{(C) \ } \frac{8}{63}\qquad \mathrm{(D) \ } \frac{1}{6}\qquad \mathrm{(E) \ } \frac{2}{7}$

 \textbf{Answer: }C



\textbf{Solution}

Let $x$ be the \href{https://artofproblemsolving.com/wiki/index.php?title=Probability}{probability} of rolling a $1$. The probabilities of rolling a $2$, $3$, $4$, $5$, and $6$ are $2x$, $3x$, $4x$, $5x$, and $6x$, respectively.

 The sum of the probabilities of rolling each number must equal 1, so

 $x+2x+3x+4x+5x+6x=1$
$21x=1$
$x=\frac{1}{21}$

 So the probabilities of rolling a $1$, $2$, $3$, $4$, $5$, and $6$ are respectively $\frac{1}{21}, \frac{2}{21}, \frac{3}{21}, \frac{4}{21}, \frac{5}{21}$, and $\frac{6}{21}$. 

 The possible combinations of two rolls that total $7$ are: $(1,6) ; (2,5) ; (3,4) ; (4,3) ; (5,2) ; (6,1)$

 The probability P of rolling a total of $7$ on the two dice is equal to the sum of the probabilities of rolling each combination. 

 $P = \frac{1}{21}\cdot\frac{6}{21}+\frac{2}{21}\cdot\frac{5}{21}+\frac{3}{21}\cdot\frac{4}{21}+\frac{4}{21}\cdot\frac{3}{21}+\frac{5}{21}\cdot\frac{2}{21}+\frac{6}{21}\cdot\frac{1}{21}=\frac{8}{63} \Rightarrow \boxed{C}$




\textbf{Solution 2 (Not as bashy)}

(Alcumus solution)On each die the probability of rolling $k$, for $1\leq k\leq 6$, is $\frac{k}{1+2+3+4+5+6}=\frac{k}{21}.$There are six ways of rolling a total of 7 on the two dice, represented by the ordered pairs $(1,6)$, $(2,5)$, $(3,4)$, $(4,3)$, $(5,2)$, and $(6,1)$. Thus the probability of rolling a total of 7 is $\frac{1\cdot6+2\cdot5+3\cdot4+4\cdot3+5\cdot2+6\cdot1}{21^2}=\frac{56}{21^2}=\boxed{\frac{8}{63}}.$




\textbf{Solution 3 (intuitive)}

There are $6$ ways to get the sum of $7$ of the dice. Let's do case by case.

 Case $1$ (Rolling a $1$ or a $6$): $\frac {1}{21} \cdot \frac {6}{21} = \frac {6}{441}$.

 Case $2$ (Rolling a $2$ or a $5$): $\frac {2}{21} \cdot \frac {5}{21} = \frac {10}{441}$. 

 Case $3$ (Rolling a $3$ or a $4$): $\frac {3}{21} \cdot \frac {4}{21} = \frac {12}{441}$. 

 The rest of the cases are symmetric to these cases above. We have $2 \cdot \frac {28}{441}$. We have $\frac {56}{441} = \frac {8}{63}$. Therefore, our answer is $\boxed {\frac {8}{63}}$




\textbf{Solution 4 (cheap)}

Notice that the ways to obtain a 7 are (6,1), (5,2), (4,3). Then, because the cases are symmetrical with (3,4), (2,5) and (6,1), look through the answer choices. You see a 4/63 and an 8/63, so obviously MAA wants you to not count the other symmetrical cases, thus giving the answer (C).

 


	\item (\textit{AMC 10 10B 2006}) Elmo makes $N$ sandwiches for a fundraiser. For each sandwich he uses $B$ globs of peanut butter at $4cent$ per glob and $J$ blobs of jam at $5cent$ per blob. The cost of the peanut butter and jam to make all the sandwiches is $\$2.53$. Assume that $B$, $J$, and $N$ are positive integers with $N>1$. What is the cost of the jam Elmo uses to make the sandwiches?


 $\mathrm{(A) \ } 1.05\qquad \mathrm{(B) \ } 1.25\qquad \mathrm{(C) \ } 1.45\qquad \mathrm{(D) \ } 1.65\qquad \mathrm{(E) \ } 1.85$

 \textbf{Answer: }D



\textbf{Solution 1}

The peanut butter and jam for each sandwich costs $4Bcent+5Jcent$, so the peanut butter and jam for $N$ sandwiches costs $N(4B+5J)cent$. 

 Setting this equal to $253cent$:

 $N(4B+5J)=253=11\cdot23$

 The only possible positive integer pairs $(N , 4B+5J)$ whose product is $253$ are: $(1,253) ; (11,23) ; (23,11) ; (253,1)$

 The first pair violates $N>1$ and the third and fourth pairs have no positive integer solutions for $B$ and $J$. 

 So, $N=11$ and $4B+5J=23$.

 The only integer solutions for $B$ and $J$ are $B=2$ and $J=3$.

 Therefore, the cost of the jam Elmo uses to make the sandwiches is $3\cdot5\cdot11=165cent$ $=  \$1.65 \Rightarrow \boxed{\mathrm{D}}$




\textbf{Solution 2}

Note as above, you get the equation $N(0.04B+0.05J)=2.53$

 Notice that we can multiply by $100$ on both sides to get whole numbers. Hence $\implies N(4B+5J)=253$

 Note that the prime factorization of $253=11\cdot23$.

 Hence, we want $4B+5J=23$ or $11$

 Now, we have two cases to test.

 Case 1: $4B+5J=23$

 Notice that we want $B\le5$ or $J\le4$

 Taking $\pmod{5}\implies 4B\equiv3\pmod{5}\implies B=2$

 Hence, $B=2,J=3$.

 Hence, the price of the jam is $3\cdot11\cdot{0.05}\implies 1.65 \implies\boxed{D}$.

 


	\item (\textit{AMC 10 10B 2006}) A triangle is partitioned into three triangles and a quadrilateral by drawing two lines from vertices to their opposite sides. The areas of the three triangles are 3, 7, and 7 as shown. What is the area of the shaded quadrilateral?


 
                    \begin{figure}[H]
                        \centering
                        \includegraphics[width=0.5\textwidth]{Problem/AMC_10_10B_2006_23_Problem_Figure_1.png}
                    \end{figure}
                    
$\mathrm{(A) \ } 15\qquad \mathrm{(B) \ } 17\qquad \mathrm{(C) \ } \frac{35}{2}\qquad \mathrm{(D) \ } 18\qquad \mathrm{(E) \ } \frac{55}{3}$

 \textbf{Answer: }D



\textbf{Solution 1}

Label the points in the figure as shown below, and draw the segment $CF$. This segment divides the quadrilateral into two triangles, let their areas be $x$ and $y$.

 
                    \begin{figure}[H]
                        \centering
                        \includegraphics[width=0.5\textwidth]{Solution/AMC_10_10B_2006_23_Solution_Figure_1.png}
                    \end{figure}
                    

 Since \href{https://artofproblemsolving.com/wiki/index.php?title=Triangle}{triangles} $AFB$ and $DFB$ share an \href{https://artofproblemsolving.com/wiki/index.php?title=Altitude}{altitude} from $B$ and have equal area, their bases must be equal, hence $AF=DF$.

 Since triangles $AFC$ and $DFC$ share an altitude from $C$ and their respective bases are equal, their areas must be equal, hence $x+3=y$. 

 Since triangles $EFA$ and $BFA$ share an altitude from $A$ and their respective areas are in the ratio $3:7$, their bases must be in the same ratio, hence $EF:FB = 3:7$.

 Since triangles $EFC$ and $BFC$ share an altitude from $C$ and their respective bases are in the ratio $3:7$, their areas must be in the same ratio, hence $x:(y+7) = 3:7$, which gives us $7x = 3(y+7)$.

 Substituting $y=x+3$ into the second equation we get $7x = 3(x+10)$, which solves to $x=\frac{15}{2}$. Then $y=x+3 = \frac{15}{2}+3 = \frac{21}{2}$, and the total area of the quadrilateral is $x+y = \frac{15}{2}+\frac{21}{2} = \boxed{\textbf{(D) }18}$.

 



\textbf{Solution 2}

Connect points $E$ and $D$. Triangles $EFA$ and $FAB$ share an altitude and their areas are in the ratio $3:7$. Their bases, $EF$ and $FB$, must be in the same $3:7$ ratio. 

 Triangles $EFD$ and $FBD$ share an altitude and their bases are in a $3:7$ ratio. Therefore, their areas are in a $3:7$ ratio and the area of triangle $EFD$ is $3$.

 Triangle $CED$ and $DEA$ share an altitude. Therefore, the ratio of their areas is equal to the ratio of bases $CE$ and $EA$. The ratio is $A:(3+3) \Rightarrow A:6$ where $A$ is the area of triangle $CED$

 Triangles $CEB$ and $EAB$ also share an altitude. The ratio of their areas is also equal to the ratio of bases $CE$ and $EA$. The ratio is $(A+3+7):(3+7) \Rightarrow (A+10):10$

 Because the two ratios are equal, we get the equation $\frac{A}{6} = \frac {A+10}{10} \Rightarrow 10A = 6A+60 \Rightarrow A = 15$. We add the area of triangle $EDF$ to get that the total area of the quadrilateral is $\boxed{\textbf{(D) }18}$.

 



\textbf{Solution 3}

We use mass points (similar to above). Let the triangle be $ABC$ with cevians (lines to opposite side) from $B$ and $C$. Let the points opposite $B$ and $C$ be $D$ and $F$ respectively and the intersection as $P$. 

 Assign masses of 1 at $B$ and $D$ since $[BPC] = [DPC]$. Then the mass at $P$ is 2. To find masses at $F$ and $C$, we let the mass at $F$ be x and the mass at $C$ be y. Then $3x = 7y$ and $y = \frac{3}{7}x$. Then $\frac{10}{7}x = 2$ since we add the masses for the fulcrum mass, and $x = \frac{7}{5}$ and $y = \frac{3}{5}$. 

 To calculate the mass at a, it is merely $\frac{7}{5} - 1 = \frac{2}{5}$ which means $\frac{[BCF]}{[ACF]} = \frac{2}{5}$ or $[ACF] = 25$. It is easy to see the answer is $\boxed{\textbf{(D) }18}$.

 


 


	\item (\textit{AMC 10 10B 2006}) Circles with centers $O$ and $P$ have radii $2$ and $4$, respectively, and are externally tangent. Points $A$ and $B$ on the circle with center $O$ and points $C$ and $D$ on the circle with center $P$ are such that $AD$ and $BC$ are common external tangents to the circles. What is the area of the concave hexagon $AOBCPD$?


 
                    \begin{figure}[H]
                        \centering
                        \includegraphics[width=0.5\textwidth]{Problem/AMC_10_10B_2006_24_Problem_Figure_1.png}
                    \end{figure}
                    
$\mathrm{(A) \ } 18\sqrt{3}\qquad \mathrm{(B) \ } 24\sqrt{2}\qquad \mathrm{(C) \ } 36\qquad \mathrm{(D) \ } 24\sqrt{3}\qquad \mathrm{(E) \ } 32\sqrt{2}$

 \textbf{Answer: }B



\textbf{Solution 1 (Similar Triangles)}

When we see this problem, it practically screams similar triangles at us. Extend $OP$ to the left until it intersects lines $AD$ and $BC$ at point $E$. Triangles $EBO$ and $ECP$ are similar, and by symmetry, so are triangles $EAO$ and $EDP$. Then, the area of kite $EAOB$ is just $4 \sqrt{2} \times 2$ and the area of kite $EDPC$ is $8 \sqrt{2} \times 4$ (using our similarity ratios). The difference of these yields the area of hexagon $AOBCPD$ or $24\sqrt{2} \Longrightarrow \boxed{\mathrm{(B)}}$.

 



\textbf{Solution 2 (Perpendiculars)}

File is too big, so go to \href{https://www.imgur.com/a/7aphGaa}{https://www.imgur.com/a/7aphGaa}

 Sorry for the wrong point names, I didn't know how to change them.

 Since a \href{https://artofproblemsolving.com/wiki/index.php?title=Tangent\_line}{tangent line} is \href{https://artofproblemsolving.com/wiki/index.php?title=Perpendicular}{perpendicular} to the \href{https://artofproblemsolving.com/wiki/index.php?title=Radius}{radius} containing the point of tangency, $\angle OAD = \angle PDA = 90^\circ$.

 Construct a perpendicular to $DP$ that goes through point $O$. Label the point of \href{https://artofproblemsolving.com/wiki/index.php?title=Intersection}{intersection} $X$.

 Clearly $OADX$ is a \href{https://artofproblemsolving.com/wiki/index.php?title=Rectangle}{rectangle}, so $DX=2$ and $PX=2$. By the \href{https://artofproblemsolving.com/wiki/index.php?title=Pythagorean\_Theorem}{Pythagorean Theorem}, $OX = \sqrt{6^2 - 2^2} = 4\sqrt{2}$.

 The area of $OADX$ is $2\cdot4\sqrt{2}=8\sqrt{2}$. The area of $OXP$ is $\frac{1}{2}\cdot2\cdot4\sqrt{2}=4\sqrt{2}$, so the area of quadrilateral $OADP$ is $8\sqrt{2}+4\sqrt{2}=12\sqrt{2}$. Using similar steps, the area of quadrilateral $OBCP$ is also $12\sqrt{2}$. Therefore, the area of hexagon $AOBCPD$ is $2\cdot12\sqrt{2}= 24\sqrt{2} \Longrightarrow \boxed{\mathrm{(B)}}$.

 
Note: Quadrilaterals $OADP$ and $OBCP$ are congruent, so they have equal areas.

 


 



\textbf{Solution 3}


                    \begin{figure}[H]
                        \centering
                        \includegraphics[width=0.5\textwidth]{Solution/AMC_10_10B_2006_24_Solution_Figure_1.png}
                    \end{figure}
                    

 Draw a line through point O parallel to $AD$ and call the point where this line meets $DP$ E. Note that $OP=6$ and $EP=2$. By the Pythagorean Theorem, $OE=4\sqrt{2}$. We created rectangle $DAOE$ and thus $AD=OE=4\sqrt{2}$. Now we know height and both top and bottom heights. The area of a trapezoid is given by: $K=\frac{b_1+b_2}{2}\cdot h$, where the $b$s are the parallel sides and $h$ is the distance between those bases. Plug stuff in to get that $[ADPO]=12\sqrt{2}$. Remember to multiply by two because there are two of these equal trapezoids, and you get $2\cdot12\sqrt{2}= 24\sqrt{2} \Longrightarrow \boxed{\mathrm{(B)}}$.

 my asymptote thingy makes the two lines meet can somebody fix that im too lazy to

 


	\item (\textit{AMC 10 10B 2006}) Mr. Jones has eight children of different ages. On a family trip his oldest child, who is 9, spots a license plate with a 4-digit number in which each of two digits appears two times. "Look, daddy!!!!" she exclaims. "That number is evenly divisible by the age of each of us kids!" "That's right," replies Mr. Jones, "and the last two digits just happen to be my age." Which of the following is \textbf{\textit{not}} the age of one of Mr. Jones's children? 


 $\mathrm{(A) \ } 4\qquad \mathrm{(B) \ } 5\qquad \mathrm{(C) \ } 6\qquad \mathrm{(D) \ } 7\qquad \mathrm{(E) \ } 8$

 \textbf{Answer: }B



\textbf{Solution 1}

Let $S$ be the set of the ages of Mr. Jones' children (in other words $i \in S$ if Mr. Jones has a child who is $i$ years old). Then $|S| = 8$ and $9 \in S$. Let $m$ be the positive integer seen on the license plate. Since at least one of $4$ or $8$ is contained in $S$, we have $4 | m$.

 We would like to prove that $5 \not\in S$, so for the sake of contradiction, let us suppose that $5 \in S$. Then $5\cdot 4 = 20 | m$ so the units digit of $m$ is $0$. Since the number has two distinct digits, each appearing twice, another digit of $m$ must be $0$. Since Mr. Jones can't be $00$ years old, the last two digits can't be $00$. Therefore $m$ must be of the form $d0d0$, where $d$ is a digit. Since $m$ is divisible by $9$, the sum of the digits of $m$ must be divisible by $9$ (see \href{https://artofproblemsolving.com/wiki/index.php?title=Divisibility\_rules/Rule\_for\_3\_and\_9\_proof}{ Divisibility rules for 9}). Hence $9 | 2d$ which implies $d = 9$. But $m = 9090$ is not divisible by $4$, contradiction. So $5 \not\in S$ and $5$ is not the age of one of Mr. Jones' kids. $\boxed{ B }$

 (We might like to check that there does, indeed, exist such a positive integer $m$. If $5$ is not an age of one of Mr. Jones' kids, then the license plate's number must be a multiple of $lcm(1,2,3,4,6,7,8,9) = 504$. Since $11\cdot504 = 5544$ and $5544$ is the only $4$ digit multiple of $504$ that fits all the conditions of the license plate's number, the license plate's number is $5544$.)

 



\textbf{Solution 2}

Alternatively, we can see that if one of Mr. Jones' children is of the age 5, then the license plate will have to end in the digit $0$ . The license plate cannot end in the digit $5$ as $2$ is a factor of the number, so it must be even. This means that the license plate would have to have two $0$ digits, and would either be of the form $XX00$ or $X0X0$ (X being the other digit in the license plate) . The condition $XX00$ is impossible as Mr. Jones can't be $00$ years old. If we separate the second condition, $X0X0$ into its prime factors, we get $X\cdot 10\cdot 101$. $101$ is prime, and therefore can't account for allowing $X0X0$ being evenly divisible by the childrens' ages.  The $10$ accounts for the 5 and one factor of 2. This leaves $X$, but no single digit number contains all the prime factors besides 5 and 2 of the childrens' ages, thus $5$ $\boxed{ B }$ can't be one of the childrens' ages.

 



\textbf{Solution 3}

Another way to do the problem is by the process of elimination. The only possible correct choices are the highest powers of each prime, $2^3=8$, $3^2=9$, $5^1=5$, and $7^1=7$, since indivisibility by any powers lower than these means indivisibility by a higher power of the prime (for example, indivisibility by $2^2=4$ means indivisibility by $2^3=8)$. Since the number is divisible by $9$, it is not the answer, and the digits have to add up to $9$. Let $(a,b)$ be the digits: since $2\cdot(a+b) | 9$, $(a+b)|9$. The only possible choices for $(a,b)$ are $(0,9)$, $(1,8)$, $(2,7)$, $(3,6)$, and $(4,5)$. The only number that works is $5544$, but it is not divisible by $5$. $\boxed{ B }$




\textbf{Solution 4}

Since 5 has the "strictest" divisibility rules of the single digit numbers, we check that case first. To be divisible by 9, the other repeating digit must be 4. Seeing that 5544 is the only arrangement of 2 4's and 2 5's that is divisible by the other single digit numbers and itself is not divisible by 5, the answer is $\boxed {(B) 5}$.-liu4505~edited by mobius247

 


\end{enumerate}